% !TEX root = 00-tutorial.tex

\chapter{Integrating Spheres}

We ignore the unscattered transmission measurements for now.

\section{Perfect measurements of total reflection/transmission (or diffuse reflection/transmission)}

Sometimes one just wants to convert a total reflection measurement to (assuming sample thickness
is 1mm and the scattering coefficient is 1mm$^{-1}$).  Here the \texttt{-V 0} option suppresses 
extra output

\begin{grayverb}
iad -r 0.5
\end{grayverb}

to get an absorption coefficient of 0.0659mm$^{-1}$.

%If the total diffuse reflection is known then use
%\begin{grayverb}
%iad -V 0 -r 0.5 -c 0
%\end{grayverb}

\section{Dual Beam Experiments}
If you happen to have a spectrophotometer with an integrating sphere that 
alternates between illuminating the sample and a reflection standard, then 
you should use the \texttt{-X} command-line and set the number of spheres to one.  

The \texttt{-X} option assumes that your measurements have already been corrected f
or most sphere effects.  

\section{Single sphere use}
Spheres are needed to collect all light reflected or transmitted by the sample.  Most
people use a single sphere; they measure sample reflectance and then move the sample
to another port to measure the transmittance.

The normalized reflection and transmission are defined six measurements 

\begin{equation}
M_R \equiv r_\mathit{std}\cdot{{R(\rdirect, \rdiffuse) - R(0, 0)\over 
                    R(r_{\mathit{std}}, r_{\mathit{std}})- R(0, 0)} }.
\label{eqmr}
\end{equation}

and

\begin{equation}
M_T \equiv {{T(\tdirect, \rdiffuse) - T_{\mathit{dark}}\over 
                    T(0, 0) - T_{\mathit{dark}}} }
\label{eqmt}
\end{equation}

where all the quantities are defined in the figure above.  Typically one collects spectra
across all wavelengths for each these six quantities.  These are collected in a \texttt{.rxt} file as below.

\begin{figure}
\begin{center}
\includegraphics[width=0.45\textwidth]{../Ch3spheresR.pdf}
\,
\includegraphics[width=0.34\textwidth]{../Ch3spheresT.pdf}
\end{center}
\caption{Measurements needed when using a single sphere.}
\end{figure}


\begin{verbatim}
IAD1  # Must be first four characters
1.38   # Index of refraction of the sample 
1.00   # Index of refraction of the top and bottom slides 
1.62   # [mm] Thickness of sample 
0.00   # [mm] Thickness of slides 
2.00   # [mm] Diameter of illumination beam 
0.99   # Reflectivity of the reflectance calibration standard

1      # Number of spheres used for the measurement

# Properties of sphere used for reflectance measurements
101.6  # [mm] Sphere Diameter
12.7   # [mm] Sample Port Diameter
12.7   # [mm] Entrance Port Diameter
0.6    # [mm] Detector Port Diameter
0.94   # Reflectivity of the sphere wall

# Properties of sphere for transmittance measurements
101.6  # [mm] Sphere Diameter
12.7   # [mm] Sample Port Diameter
12.7   # [mm] Third Port Diameter
0.6    # [mm] Detector Port Diameter
0.94   # Reflectivity of the sphere wall

L      r         t     g
852 0.315087 0.606721 0.9
853 0.333245 0.590481 0.9

\end{verbatim}

By creating such a file, all the information about a particular experiment is captured in one place.  If this file is called
\texttt{example.rxt} then 

\begin{grayverb}
iad example.rxt
\end{grayverb}
